\scanpage{235}
\mspage{185}
% L1–2
\[
H^n(X,F_{\unsure{X|x}})\unsure{-|\to}\mathrm{Hom}_\pi(\pi_n(Z,x),\underset{\substack{\Vert\\ F_y}}{M})
\]
% L3–4
n'est autre que la composée de la flèche précédente (page précédente) pour $\unsure{G|C}=F_X$, avec
% L5
\[
\mathrm{Hom}_\pi(\pi_n(X,x),M)\to\mathrm{Hom}_\pi(\pi_n(Z,x),M)
\]
% L6
déduite de
% L7
\[
\pi_n(Z,x)\unsure{-|\to}\pi_n(X,x).
\]
% L8–13
Or cette flèche est nulle, par l'hypothèse $b_n$. Donc la flèche $\alpha$ dans la suite exacte page 181 est nulle, donc $H^n(Y,F)\to H^n(X,F_X)$ est iso, qed (mais il y a pas mal de travail de vérifications à faire\ldots)

% L14–19
$c_n\Rightarrow a_{n-1}$. On peut procéder par récurrence sur $n$ (NB l'énoncé a un sens même pour $n=1$, où il est prouvé, et $\unsure{=}$ $c_1\Rightarrow b_1$, et $b_1\Rightarrow a_0$, \inserted{plus généralement $b_n\Rightarrow a_{n-1}$,} immédiat par suite exacte d'homotopie.)
% L20–23
Donc on sait déjà, en plus de l'hypothèse $c_n$, qu'on a $a_{n-2}$ i.e. $Z$ est $(n-\unsure{2|3})$-connexe. Il faut prouver qu'il est \inserted{$\unsure{=}$} $(n-1)$-connexe, i.e. $\pi_{n-1}(Z)=0$. Le cas $n=2$ demande
% archive sheet number: 235
